\section{Related Work} \label{sec:2}

\subsection{Current Diagnostic Practice and the Case for Objective Biomarkers} \label{sec:2.1}
Clinical diagnosis of PD remains anchored to the MDS Clinical Diagnostic Criteria \cite{Postuma:2015}: a judgment, made by a trained examiner, of whether bradykinesia is present together with rest tremor and/or rigidity, refined by supportive features and exclusion criteria and typically corroborated by the patient's response to levodopa therapy. Because this judgment rests on the visual and manual assessment of motor signs that can be subtle, intermittent, or altogether absent early in the disease course, diagnostic accuracy is known to vary considerably with examiner experience and disease duration, and misdiagnosis remains common enough at early stages to be a recurring concern in the movement-disorders literature \cite{Rizzo:2016}. This has motivated a broader search for objective biomarkers that do not depend on a clinician's subjective interpretation of motor signs. Dopaminergic imaging (DAT-SPECT, PET), cerebrospinal-fluid and blood-based $\alpha$-synuclein assays, genetic testing, and, more recently, wearable and other digital-sensing modalities have all been investigated as candidate objective markers \cite{Sun:2024}; each can support diagnosis or staging, but each also carries some combination of high cost, invasiveness, specialized equipment, or limited accessibility outside tertiary care centers \cite{Sun:2024}. Voice and speech recordings, by contrast, can be acquired with nothing more than a standard microphone, at negligible marginal cost, and without any physical contact with the patient, a combination of properties that renders them an unusually attractive candidate among the emerging digital-biomarker modalities, and the specific focus of this study.

\subsection{Acoustic and Speech Biomarkers of Parkinsonian Dysphonia} \label{sec:2.2}
Hypokinetic dysarthria, characterized by reduced vocal loudness, monopitch, imprecise articulation, and breathy or hoarse voice quality, is estimated to affect the large majority of individuals with PD at some point in the disease course \cite{Ho:1999, Cao:2025}. Its physiological basis lies in rigidity and bradykinesia affecting the laryngeal, respiratory, and articulatory musculature: reduced vocal-fold adduction force and control lead to incomplete or irregular glottal closure, reduced respiratory support lowers subglottal pressure and vocal intensity, and reduced range of motion in the articulators blurs consonant and vowel targets \cite{Yang:2020}. These physiological changes are reflected in several well-established acoustic correlates, namely increased cycle-to-cycle perturbation in fundamental frequency (jitter) and amplitude (shimmer), reduced harmonics-to-noise ratio (HNR) reflecting incomplete glottal closure, reduced vowel space area reflecting imprecise articulation, and altered nonlinear dynamical measures of vocal-fold vibration, including recurrence period density entropy (RPDE), detrended fluctuation analysis (DFA), and pitch period entropy (PPE), that are more sensitive than classical perturbation measures to the aperiodic, weakly chaotic component of parkinsonian phonation \cite{Yang:2020}. Given that these physiological changes can begin before overt motor signs are clinically apparent, acoustic and broader speech biomarkers have increasingly been proposed as tools for prodromal detection and longitudinal monitoring, not merely post-diagnosis screening \cite{Rusz:2024, Cao:2025}. This physiological grounding is what would make the specific acoustic features a wrapper-based pipeline selects clinically meaningful, not merely predictive; Section \ref{sec:5} returns to this point as a direction for future work, since the present study's logged results support comparing how many features are selected but not which ones.

\subsection{Metaheuristic Feature Selection for PD Detection: Progress and Methodological Gaps} \label{sec:2.3}
Wrapper-based feature selection couples a search procedure with a downstream classifier's cross-validated performance as the objective to optimize. Population-based metaheuristics, being derivative-free and requiring no assumption of a smooth or convex fitness landscape, have become a standard choice of search procedure for this task across biomedical applications generally \cite{BarreraGarcia:2023}. Within voice-based PD detection specifically, a fast-growing line of work applies swarm- and evolution-inspired algorithms to select acoustic feature subsets ahead of a downstream classifier, most commonly encoding the search space as a binary mask over the candidate feature set: particle-swarm-, butterfly-, and whale-inspired variants \cite{Hashemi:2026}, particle swarm optimization jointly tuning the feature subset and classifier hyperparameters \cite{Eliwa:2025}, a modified grey wolf optimizer \cite{Santhosh:2025}, a modified hunger games search \cite{Hashim:2023}, a quantum mayfly optimizer \cite{Mansour:2024}, a genetic algorithm with filter-based pre-selection \cite{Ali:2023}, a beluga whale optimizer \cite{Ganesan:2025}, and hybrids chaining two algorithms in sequence, namely a whale-to-grey-wolf cascade \cite{AlNajjar:2024} and a chaotic grey wolf-dragonfly cascade \cite{Appati:2026}.

As discussed in Section \ref{sec:1}, the classifier consuming the selected features is fixed and predetermined in the large majority of this work, rarely optimized jointly with the feature subset within a single search process. Even where classifier hyperparameters are tuned alongside the feature subset \cite{Eliwa:2025}, the classifier type itself is still chosen in advance rather than searched over. This omission is not merely a matter of scope: because different classifiers weight, combine, and interact with input features in fundamentally different ways, the feature subset best suited to one classifier need not be the one best suited to another, so fixing the classifier a priori risks discarding feature subsets that would perform poorly under the assumed classifier but well under an alternative one that was never considered. Extending the search to include the classifier itself, however, is not a trivial enlargement of the existing formulations; it substantially increases the dimensionality and, more importantly, the combinatorial complexity of the search space, since feature and classifier choices interact rather than contribute independently to fitness, which is precisely the setting in which premature convergence and loss of population diversity become most damaging to search quality. This is the gap that the encoding proposed in this study is designed to close, by searching a higher-dimensional, joint feature-and-classifier space without sacrificing search quality relative to feature-selection-only formulations.
